\documentclass{article}
\usepackage{relsize}
\usepackage{amssymb}
\usepackage{geometry}

% Reduce margins
\geometry{
    a4paper,
    margin=1in % You can adjust this value (e.g., 0.75in, 2cm, etc.)
}

\newcommand{\subtitlerelsize}{1} %relative size: integer value
\newcommand{\subtitlelinesep}{0.2em} %line separation: a LaTeX length

\title{Evaluationsbogen Agenten-Frameworks}
\date{}

\begin{document}
    
\maketitle
    
\vspace{-1.5cm}

\begin{center}
    \textbf{Framework:} AutoGen
    \textbf{Testzeitraum:} x Tage
\end{center}

\section*{Teil A}

\newcommand{\evalquestion}[2]{
    #1
    \begin{flushleft}
        \ifnum#2=1 $\boxtimes$ Stimme gar nicht zu \else $\square$ Stimme gar nicht zu \fi \quad
        \ifnum#2=2 $\boxtimes$ Eher nicht \else $\square$ Eher nicht \fi \quad
        \ifnum#2=3 $\boxtimes$ Neutral \else $\square$ Neutral \fi \quad
        \ifnum#2=4 $\boxtimes$ Eher ja \else $\square$ Eher ja \fi \quad
        \ifnum#2=5 $\boxtimes$ Stimme voll zu \else $\square$ Stimme voll zu \fi
    \end{flushleft}
}

\noindent
\evalquestion{1. Ich denke, dass ich dieses Framework häufiger verwenden möchte.}{0}

\noindent
\evalquestion{2. Ich fand das Framework unnötig komplex.}{0}

\noindent
\evalquestion{3. Ich dachte, das Framework war einfach zu bedienen.}{0}

\noindent
\evalquestion{4. Ich denke, dass ich die Unterstützung eines technischen Supports brauche, um dieses Framework nutzen zu können.}{0}

\noindent
\evalquestion{5. Ich fand, die verschiedenen Funktionen in diesem Framework waren gut integriert.}{0}

\noindent
\evalquestion{6. Ich dachte, dass dieses Framework nicht konsistent genug war.}{0}

\noindent
\evalquestion{7. Ich würde mir vorstellen, dass die meisten Informatiker sehr schnell lernen würden, dieses Framework zu benutzen.}{0}

\noindent
\evalquestion{8. Ich fand dieses Framework sehr umständlich zu benutzen.}{0}

\noindent
\evalquestion{9. Ich habe mich sehr selbstsicher gefühlt, dieses Framework zu verwenden.}{0}

\noindent
\evalquestion{10. Ich musste eine Menge lernen, bevor ich mit diesem Framework loslegen konnte.}{0}

\section*{Teil B}

\noindent
\evalquestion{1. Die technische Dokumentation des Frameworks ist vollständig, aktuell und verständlich.}{0}

\noindent
2. Welche wichtigen Informationen waren in der Dokumentation nicht oder nur schwer
auffindbar? \\

Die allermeisten Informationen waren dokumentiert und der Quick Start war verständlich. 
Im Vergleich mit anderen Frameworks (insbesondere OpenAI SDK) war aber dennoch das differnzierte 
Logging ein Problem und wurde von der Dokumentation nicht ausreichend erklärt. Um einige Events 
auszuschließen musste die konkrete API konsultiert werden um die Event-Typen zu finden 
(genauer in Workarounds). In Anbetracht der Tatsache, dass das Logging besonders für Einsteiger
ein guter Ansatzpunkt für das Debugging und für die Nachvollziehbarkeit ist, war das ärgerlich. 
% TODO: Bei OpenAI war die EInbindung von anderen LLMs und der ResponseModelle sehr schwer zu finden. 
% In Work-Arounds: Warum Response Modelle?

Des weiteren steckte ich anfangs lange bei den vorgefertigten Tools fest (namentlich dem FileServerTool),
welches meiner Meinung nach für das Beschaffen und/oder Lesen der Dateien geeignet war, was allerdings
nicht der Fall war. Leider war der konkrete Anwendungsfall nicht in der Dokumentation beschrieben. \\

\noindent
3. Gab es Aspekte, die sich im praktischen Einsatz als besonders durchdacht oder
problematisch erwiesen haben – unabhängig von der Dokumentation? \\

\noindent
4. Welche Lösungen oder Workarounds mussten implementiert werden, um bestimmte
Anforderungen oder Anwendungsfälle umzusetzen? \\


\noindent
5. Gab es Funktionen oder Eigenschaften des Frameworks, die Sie positiv oder negativ
überrascht haben? \\
Vergleichend mit den anderen Frameworks auf alle Fälle 

\noindent
\evalquestion{6. Wie gut lässt sich das von Ihnen mit dem Framework erstellte System (auch in kleinerem
Maßstab) im Nachhinein ändern oder erweitern?}{0}
Besonders durch die RoundRobin-Kommunikation für Teams ist eine Erweiterung sehr gut realisierbar, da
sich um die Koordination nicht gekümmert werden muss, - zumindest im Anwendungsfall dieser Experimente - 
da die Koordination autonom erfolgt. Mögliche Änderungen an dem Task-Problem, beispielsweise, müssten meiner
Ansicht nach in jedem Framework getätigt werden, so dass das kein negativen Punkt darstellt. \\


\noindent
\evalquestion{7. Wie gut wäre das von Ihnen erstellte System mit vertretbarem Aufwand in einem
produktiven Szenario einsetzbar?}{0}

\noindent
8. Wie haben Sie das System hauptsächlich erstellt?
\begin{flushleft}
    $\square$ Überwiegend durch Programmierung \quad
    $\square$ Über Konfiguration (z. B. YAML, JSON) \quad
    $\square$ Über grafische Benutzeroberfläche (GUI) \quad
\end{flushleft}

\noindent
9. Wie schätzen Sie den Aufwand und die Flexibilität dieser Art der Systemerstellung
( siehe Frage 8 ) ein?



\end{document}